% GENERATED FILE. Edit canonical lessons, then run npm run generate:books.
% canonical-source-hash: fnv1a64:8c5289fb29a3ac72
% canonical-lessons: ES-C735-esqueleto, ES-C735-pulmon, ES-C735-estomago, ES-C735-cerebro, ES-C735-latir, ES-R735-repaso-728-735

\chapter{Por dentro del cuerpo}
\label{ch:es-por-dentro-del-cuerpo}
\hlchaptermodality{735}

\begin{chapteropening}
\textbf{By the end of this chapter:} \emph{I can name the skeleton and inner organs and say that my heart is beating fast.}

Complete the last lesson of chapter 735: I can name the skeleton and inner organs and say that my heart is beating fast.
\end{chapteropening}

\section[el esqueleto]{el esqueleto — the skeleton}

\label{lesson:ES-C735-esqueleto}



\begin{quote}
Before the new one: say the Spanish for to bleed, then the Spanish for the scar.
\end{quote}

\subsection*{You'll want to know: el esqueleto}
\textbf{el esqueleto} — \textquotedblleft{}the skeleton\textquotedblright{}.

\begin{quote}
El esqueleto humano tiene 206 huesos.
\end{quote}

\begin{quote}
\emph{The human skeleton has 206 bones.}
\end{quote}

\begin{cousinweb}
\emph{esqueleto} comes from Greek \emph{skeletós} \textquotedblleft{}dried up, a dried body\textquotedblright{} (EtymDB), from \emph{skéllein} \textquotedblleft{}to dry up\textquotedblright{}.
\end{cousinweb}

\begin{grammarlens}[title={What you've built}]
One more word for telling, explaining and describing at B1.
\end{grammarlens}

\subsection*{Your turn}
\begin{itemize}
\raggedright
  \item \emph{Say it:} \emph{el esqueleto}
  \item \emph{Say it:} \emph{el esqueleto}, once more
  \item \emph{Say it:} a sentence of your own with \emph{el esqueleto}
  \item \emph{Recall:} say the Spanish for to bleed, then the Spanish for the scar, then say \emph{el esqueleto} again
\end{itemize}

\begin{tcolorbox}[breakable,colback=teal!4,colframe=teal!35!black,title={Before you move on}]
What does el esqueleto mean? (\textquotedblleft{}The skeleton\textquotedblright{}.) Say it once more.
\end{tcolorbox}
\section[el pulmón]{el pulmón — the lung}

\label{lesson:ES-C735-pulmon}



\begin{quote}
Before the new one: say the Spanish for the scar, then the Spanish for the skeleton.
\end{quote}

\subsection*{You'll want to know: el pulmón}
\textbf{el pulmón} — \textquotedblleft{}the lung\textquotedblright{}.

\begin{quote}
Fumar es malo para los pulmones.
\end{quote}

\begin{quote}
\emph{Smoking is bad for the lungs.}
\end{quote}

Plural \emph{los pulmones}.

\begin{cousinweb}
\emph{pulmón} comes from Latin \emph{pulmō}, \emph{pulmōnem} \textquotedblleft{}lung\textquotedblright{} (EtymDB).
\end{cousinweb}

\begin{grammarlens}[title={What you've built}]
One more word for telling, explaining and describing at B1.
\end{grammarlens}

\subsection*{Your turn}
\begin{itemize}
\raggedright
  \item \emph{Say it:} \emph{el pulmón}
  \item \emph{Say it:} \emph{el pulmón}, once more
  \item \emph{Say it:} a sentence of your own with \emph{el pulmón}
  \item \emph{Recall:} say the Spanish for the scar, then the Spanish for the skeleton, then say \emph{el pulmón} again
\end{itemize}

\begin{tcolorbox}[breakable,colback=teal!4,colframe=teal!35!black,title={Before you move on}]
What does el pulmón mean? (\textquotedblleft{}The lung\textquotedblright{}.) Say it once more.
\end{tcolorbox}
\section[el estómago]{el estómago — the stomach}

\label{lesson:ES-C735-estomago}



\begin{quote}
Before the new one: say the Spanish for the skeleton, then the Spanish for the lung.
\end{quote}

\subsection*{You'll want to know: el estómago}
\textbf{el estómago} — \textquotedblleft{}the stomach\textquotedblright{}.

\begin{quote}
Me duele el estómago.
\end{quote}

\begin{quote}
\emph{My stomach hurts.}
\end{quote}

\emph{Dolor de estómago} — stomach ache.

\begin{cousinweb}
\emph{estómago} is a learned borrowing from Latin \emph{stomachus}, which came from Greek \emph{stómakhos}, built on \emph{stóma} \textquotedblleft{}mouth\textquotedblright{} (EtymDB).
\end{cousinweb}

\begin{grammarlens}[title={What you've built}]
One more word for telling, explaining and describing at B1.
\end{grammarlens}

\subsection*{Your turn}
\begin{itemize}
\raggedright
  \item \emph{Say it:} \emph{el estómago}
  \item \emph{Say it:} \emph{el estómago}, once more
  \item \emph{Say it:} a sentence of your own with \emph{el estómago}
  \item \emph{Recall:} say the Spanish for the skeleton, then the Spanish for the lung, then say \emph{el estómago} again
\end{itemize}

\begin{tcolorbox}[breakable,colback=teal!4,colframe=teal!35!black,title={Before you move on}]
What does el estómago mean? (\textquotedblleft{}The stomach\textquotedblright{}.) Say it once more.
\end{tcolorbox}
\section[el cerebro]{el cerebro — the brain}

\label{lesson:ES-C735-cerebro}



\begin{quote}
Before the new one: say the Spanish for the lung, then the Spanish for the stomach.
\end{quote}

\subsection*{You'll want to know: el cerebro}
\textbf{el cerebro} — \textquotedblleft{}the brain\textquotedblright{}.

\begin{quote}
El cerebro necesita dormir.
\end{quote}

\begin{quote}
\emph{The brain needs sleep.}
\end{quote}

\emph{El cerebro del grupo} — the brains of the group.

\begin{cousinweb}
\emph{cerebro} is usually described as a learned borrowing from Latin \emph{cerebrum} \textquotedblleft{}brain\textquotedblright{}.
\end{cousinweb}

\begin{grammarlens}[title={What you've built}]
One more word for telling, explaining and describing at B1.
\end{grammarlens}

\subsection*{Your turn}
\begin{itemize}
\raggedright
  \item \emph{Say it:} \emph{el cerebro}
  \item \emph{Say it:} \emph{el cerebro}, once more
  \item \emph{Say it:} a sentence of your own with \emph{el cerebro}
  \item \emph{Recall:} say the Spanish for the lung, then the Spanish for the stomach, then say \emph{el cerebro} again
\end{itemize}

\begin{tcolorbox}[breakable,colback=teal!4,colframe=teal!35!black,title={Before you move on}]
What does el cerebro mean? (\textquotedblleft{}The brain\textquotedblright{}.) Say it once more.
\end{tcolorbox}
\section[latir]{latir — to beat (of the heart)}

\label{lesson:ES-C735-latir}



\begin{quote}
Before the new one: say the Spanish for the stomach, then the Spanish for the brain.
\end{quote}

\subsection*{You'll want to know: latir}
\textbf{latir} — \textquotedblleft{}to beat (of the heart)\textquotedblright{}.

\begin{quote}
Me late muy rápido el corazón.
\end{quote}

\begin{quote}
\emph{My heart is beating very fast.}
\end{quote}

\emph{El latido} is the heartbeat.

\begin{cousinweb}
\emph{latir} comes from Medieval Latin \emph{glattīre} \textquotedblleft{}to yelp\textquotedblright{} (EtymDB). Old \emph{latir} was a dog's yelping; the steady, repeated sound gave the heartbeat.
\end{cousinweb}

\begin{grammarlens}[title={What you've built}]
One more word for telling, explaining and describing at B1.
\end{grammarlens}

\subsection*{Your turn}
\begin{itemize}
\raggedright
  \item \emph{Say it:} \emph{latir}
  \item \emph{Say it:} \emph{latir}, once more
  \item \emph{Say it:} a sentence of your own with \emph{latir}
  \item \emph{Recall:} say the Spanish for the stomach, then the Spanish for the brain, then say \emph{latir} again
\end{itemize}

\begin{tcolorbox}[breakable,colback=teal!4,colframe=teal!35!black,title={Before you move on}]
What does latir mean? (\textquotedblleft{}To beat (of the heart)\textquotedblright{}.) Say it once more.
\end{tcolorbox}
\section[(dialogue)]{Repaso de los capítulos 728–735}

\label{lesson:ES-R735-repaso-728-735}



\begin{quote}
Say the Spanish for wise and wisdom.
\end{quote}

\subsection*{Your turn}
Say these aloud:

\begin{itemize}
\raggedright
  \item wise, wisdom, clever, skilful, wit
  \item stubborn, lazy, laziness, nosy, curiosity
  \item charming, charm, attractive, beauty, elegance
  \item mature, maturity, to mature, optimistic, pessimistic
  \item healthy, to look after oneself, illness, to fall ill, to get worse
  \item the flu, the cold, having a cold, to infect, to sneeze
  \item to burn oneself, the burn, the bandage, to bleed, the scar
  \item the skeleton, the lung, the stomach, the brain, to beat
\end{itemize}

\begin{tcolorbox}[breakable,colback=teal!4,colframe=teal!35!black,title={Before you move on}]
Wise? (\textbf{sabio}.) Wisdom? (\textbf{la sabiduría}.) Clever? (\textbf{astuto}.) Skilful? (\textbf{hábil}.) Wit? (\textbf{el ingenio}.) Stubborn? (\textbf{terco}.) Lazy? (\textbf{perezoso}.) Laziness? (\textbf{la pereza}.) Nosy? (\textbf{entrometido}.) Curiosity? (\textbf{la curiosidad}.) Charming? (\textbf{encantador}.) Charm? (\textbf{el encanto}.) Attractive? (\textbf{atractivo}.) Beauty? (\textbf{la belleza}.) Elegance? (\textbf{la elegancia}.) Mature? (\textbf{maduro}.) Maturity? (\textbf{la madurez}.) To mature? (\textbf{madurar}.) Optimistic? (\textbf{optimista}.) Pessimistic? (\textbf{pesimista}.) Healthy? (\textbf{saludable}.) To look after oneself? (\textbf{cuidarse}.) Illness? (\textbf{la enfermedad}.) To fall ill? (\textbf{enfermarse}.) To get worse? (\textbf{empeorar}.) The flu? (\textbf{la gripe}.) The cold? (\textbf{el resfriado}.) Having a cold? (\textbf{constipado}.) To infect? (\textbf{contagiar}.) To sneeze? (\textbf{estornudar}.) To burn oneself? (\textbf{quemarse}.) The burn? (\textbf{la quemadura}.) The bandage? (\textbf{la venda}.) To bleed? (\textbf{sangrar}.) The scar? (\textbf{la cicatriz}.) The skeleton? (\textbf{el esqueleto}.) The lung? (\textbf{el pulmón}.) The stomach? (\textbf{el estómago}.) The brain? (\textbf{el cerebro}.) To beat? (\textbf{latir}.)
\end{tcolorbox}
